\section{复用表格样式与列定义}
\key{\gradtablestyle{名称}{表格选项}} 保存整张表的外观，
\key{\gradtableschema{名称}{列声明}} 保存列顺序、标题、类型和列参数。
在环境选项中使用 \key{style=名称}，在环境内部使用 \key{\graduseschema{名称}}。
定义一般放在导言区；重复定义同名项目会替换之前的定义。

\begin{example}[text and listing,sidebyside=false]{示例 64：复用同一套样式与列定义}
\gradtablestyle{benchmark}{rules=booktabs,header align=c,
  text width=30mm,graphic width=40mm,label width=12mm,
  row stretch=1.3,group color=gradbarsBlue!8}
\gradtableschema{benchmark}{
  \gradcolumn{method}{方法}
  \gradcolumn[type=bar,mark=both,
    options={max=100,precision=1}]{score}{得分}
}
\begin{gradtable}[style=benchmark,header align=c]
  \graduseschema{benchmark}
  \gradgroup{验证集}
  \gradrow{方法 A,81.5}
  \gradrow{方法 B,88.2}
\end{gradtable}

\begin{gradtable}[style=benchmark,rules=grid,header align=c]
  \graduseschema{benchmark}
  \gradgroup{测试集}
  \gradrow{方法 A,79.6}
  \gradrow{方法 B,85.3}
\end{gradtable}
\end{example}
选项按从左到右的顺序执行；放在 style 后面的参数覆盖命名样式。
样式可以引用其他已定义样式，列定义也可以组合其他列定义；循环引用会报告错误。
列标识符在一张表内必须唯一，列定义必须在第一个数据行或分组标题之前调用。
每张表分别计算列宽和最优值，不会把前一张表的数据带入下一张表。
默认 \key{gradtable} 不分页；设置 \key{long=true} 可启用跨页模式，详见后文。默认不生成汇总行，需显式设置 \key{aggregate} 并声明 \key{\gradsummary}。
